\documentclass[10pt]{article}

% COMP 2300 Induction Core companion: student-facing reference and practice.
\usepackage[letterpaper,margin=0.56in,includehead,includefoot,
  headheight=15pt,headsep=7pt,footskip=20pt]{geometry}
\usepackage[T1]{fontenc}
\usepackage{amsmath,amssymb,mathtools}
\usepackage{microtype}
\usepackage[dvipsnames,svgnames]{xcolor}
\usepackage{multicol}
\usepackage{enumitem}
\usepackage{tabularx,booktabs,array}
\usepackage[most]{tcolorbox}
\usepackage{fancyhdr}
\usepackage{titlesec}
\usepackage[hidelinks]{hyperref}

\definecolor{CoreNavy}{HTML}{18364A}
\definecolor{CoreTeal}{HTML}{176B6B}
\definecolor{CoreGold}{HTML}{9A6500}
\definecolor{CoreRed}{HTML}{963D38}
\definecolor{CoreBlueTint}{HTML}{EDF4F8}
\definecolor{CoreTealTint}{HTML}{EDF8F6}
\definecolor{CoreGoldTint}{HTML}{FFF7E5}
\definecolor{CoreRedTint}{HTML}{FCEFED}

\setlength{\parindent}{0pt}
\setlength{\parskip}{2.5pt}
\setlength{\columnsep}{0.25in}
\setlength{\columnseprule}{0.25pt}
\setlength{\emergencystretch}{1.2em}
\setlength{\abovedisplayskip}{3pt}
\setlength{\belowdisplayskip}{3pt}
\setlist{leftmargin=*,itemsep=1pt,topsep=1.5pt,parsep=0pt,partopsep=0pt}
\raggedbottom

\titleformat{\section}{\Large\bfseries\color{CoreNavy}}{}{0pt}{}
  [\vspace{-3pt}\color{CoreNavy}\rule{\linewidth}{0.7pt}]
\titlespacing*{\section}{0pt}{3pt}{4pt}

\pagestyle{fancy}
\fancyhf{}
\fancyhead[L]{\textbf{COMP 2300}}
\fancyhead[C]{\textbf{Induction Core Companion}}
\fancyhead[R]{Definitions and Worked Examples}
\fancyfoot[L]{Compact Core sections A--C}
\fancyfoot[C]{\thepage}
\fancyfoot[R]{Study reference---not used during a Core}
\renewcommand{\headrulewidth}{0.5pt}
\renewcommand{\footrulewidth}{0.3pt}

\tcbset{enhanced,before skip=2.5pt,after skip=2.5pt,boxrule=0.55pt,arc=1.2mm,
  left=4pt,right=4pt,top=3pt,bottom=3pt,fonttitle=\bfseries}
\newtcolorbox{definitionbox}[1]{title={#1},colback=CoreBlueTint,
  colframe=CoreNavy,colbacktitle=CoreNavy,coltitle=white}
\newtcolorbox{methodbox}[1]{title={#1},colback=CoreTealTint,
  colframe=CoreTeal,colbacktitle=CoreTeal,coltitle=white}
\newtcolorbox{workedbox}[1]{title={#1},colback=white,
  colframe=CoreNavy,colbacktitle=CoreBlueTint,coltitle=CoreNavy}
\newtcolorbox{warningbox}[1]{title={#1},colback=CoreRedTint,
  colframe=CoreRed,colbacktitle=CoreRed,coltitle=white}
\newtcolorbox{checklistbox}[1]{title={#1},colback=CoreGoldTint,
  colframe=CoreGold,colbacktitle=CoreGold,coltitle=white}
\newcommand{\term}[1]{\textbf{\color{CoreNavy}#1}}

\begin{document}
\small

\section*{A. Build a Valid Induction Plan}
\begin{multicols}{2}
\begin{definitionbox}{The five required pieces}
An induction proof establishes a proposition $P(n)$ for every integer
$n\ge n_0$ by giving:
\begin{enumerate}
  \item a precise statement of $P(n)$ and its domain;
  \item a verified \term{base case} $P(n_0)$;
  \item an \term{induction hypothesis} for an arbitrary $k\ge n_0$;
  \item an \term{induction step} deriving $P(k+1)$ from the hypothesis; and
  \item a conclusion invoking induction.
\end{enumerate}
The hypothesis is temporary: it may be used only while proving the next case.
\end{definitionbox}

\begin{definitionbox}{Ordinary versus strong induction}
\term{Ordinary induction} assumes only $P(k)$ and proves $P(k+1)$.
\term{Strong induction} assumes every case
\[
P(n_0),P(n_0+1),\ldots,P(k)
\]
and proves $P(k+1)$. The two principles are logically equivalent; choose the
form that naturally supplies the earlier cases needed in the step.
\end{definitionbox}

\begin{methodbox}{A planning template}
Before calculating, write:
\begin{enumerate}
  \item $P(n)$: the exact claim with the variable visible;
  \item Base: substitute the first permitted value;
  \item IH: ``Fix arbitrary $k\ge n_0$ and assume \ldots'';
  \item Goal: write $P(k+1)$ explicitly;
  \item Bridge: identify the expression where the IH will enter.
\end{enumerate}
For strong induction, state the entire range of assumed cases.
\end{methodbox}

\columnbreak

\begin{workedbox}{A.1: prime factorization setup}
Claim: every integer $n\ge2$ is a product of one or more primes.

\textbf{Base.} $P(2)$ holds because $2$ is prime.

\textbf{Strong IH.} Fix arbitrary $k\ge2$ and assume $P(j)$ holds for every
integer $j$ with $2\le j\le k$.

\textbf{Goal.} Prove $P(k+1)$: that $k+1$ is a product of primes.

\textbf{Cases.} If $k+1$ is prime, it is already a product of one prime. If
it is composite, write $k+1=ab$ with $2\le a,b\le k$; the strong IH applies
to both factors. These are the two cases the short Core problem asks you to
identify; do not assume a factorization before entering the composite case.
\end{workedbox}

\begin{workedbox}{A similar ordinary-induction plan}
\textbf{1. State and start.} For
$P(n):1+3+\cdots+(2n-1)=n^2$, use base $n=1$.
\textbf{2. Assume the $k$ case:}
\[
1+3+\cdots+(2k-1)=k^2.
\]
\textbf{3. Write the target.} The next new term is
$2(k+1)-1=2k+1$. \textbf{4. Build the bridge:}
\[
k^2+(2k+1)=(k+1)^2.
\]
The equality $k^2=1+\cdots+(2k-1)$ is the exact point where the hypothesis
enters; then the display reaches the $k+1$ formula.
\end{workedbox}

\begin{warningbox}{Frequent A errors}
Do not write ``assume the claim is true'' without naming the case(s). Do not
assume $P(k+1)$. For strong induction, do not omit the lower and upper limits
of the assumed range. Verify every needed initial case for a multi-step
recurrence.
\end{warningbox}

\begin{checklistbox}{A readiness}
Can I state the base, hypothesis, next-case goal, and useful case split without
silently assuming the conclusion?
\end{checklistbox}
\end{multicols}

\clearpage
\section*{B. Diagnose and Repair an Induction Step}
\begin{multicols}{2}
\begin{definitionbox}{Circular reasoning}
An induction step must begin from the induction hypothesis $P(k)$ (or the
strong range through $k$) and derive $P(k+1)$. Beginning with $P(k+1)$ and
manipulating it is circular unless every step is a reversible equivalence and
the proof is explicitly reorganized into a forward derivation.
\end{definitionbox}

\begin{methodbox}{A repair workflow}
\begin{enumerate}
  \item Write the actual $P(k)$ and $P(k+1)$ side by side.
  \item Delete any use of the target $P(k+1)$ as a premise.
  \item Start with the left side or defining object in $P(k+1)$.
  \item Isolate a copy of the $k$-case and cite the IH exactly there.
  \item Finish the algebra or inequality, then state the conclusion.
\end{enumerate}
\end{methodbox}

\begin{workedbox}{B.1: repair the exponential inequality}
Claim: $2^n\ge n+1$ for every $n\ge0$. The faulty step assumes
$2^{k+1}\ge k+2$, which is exactly the target. \textbf{Step 1: delete that
assumption.}

\textbf{Step 2: use the correct hypothesis} $2^k\ge k+1$.
\textbf{Step 3: start from the next-case expression and move forward:}
\begin{align*}
2^{k+1}&=2\cdot2^k\\
       &\ge2(k+1) &&\text{by the IH}\\
       &\ge k+2 &&\text{because }k\ge0.
\end{align*}
\textbf{Step 4: close.} Thus $P(k+1)$ follows. Notice that the IH is used once and its direction is
preserved when multiplying by the positive number $2$.
\end{workedbox}

\columnbreak

\begin{definitionbox}{Other structural failures}
\begin{itemize}
  \item \term{Missing base case}: the step alone starts no chain.
  \item \term{Too few bases}: a recurrence using two earlier terms needs enough
  initial cases.
  \item \term{Wrong next term}: the added term must be computed at $k+1$.
  \item \term{Unjustified strengthening}: proving a stronger statement is fine,
  but it needs its own base and step.
  \item \term{Domain drift}: divisibility, inequalities, or recurrences may
  require a restricted integer domain.
\end{itemize}
\end{definitionbox}

\begin{workedbox}{Repair a wrong next term}
Suppose $P(n)$ is $1+3+\cdots+(2n-1)=n^2$. After assuming $P(k)$, the next
term is $2(k+1)-1=2k+1$, not $2k-1$. Hence
\[
1+3+\cdots+(2k-1)+(2k+1)
=k^2+2k+1=(k+1)^2.
\]
The index, not visual pattern-matching, determines the new term.
\end{workedbox}

\begin{workedbox}{Two-step recurrences need two bases}
If $a_n=a_{n-1}+a_{n-2}$ and a property of $a_n$ is proved using the two
previous cases, verify both initial values. In a strong-induction version,
the hypothesis must include enough earlier indices for the recurrence to be
defined. One base case cannot support the first step that needs two values.
\end{workedbox}

\begin{warningbox}{Invalid shortcuts}
Do not ``cancel'' an inequality without checking signs. Do not replace the
induction hypothesis with the desired conclusion. Do not test several values
and call that an induction proof. A pattern suggests a claim; the proof must
cover every allowed integer.
\end{warningbox}

\begin{checklistbox}{B readiness}
Can I point to the first circular or unsupported line, state the correct IH,
and replace the broken step with a forward chain that explicitly cites it?
\end{checklistbox}
\end{multicols}

\clearpage
\section*{C. The Essential Binomial-Sum Proof}
\begin{multicols}{2}
\begin{definitionbox}{Binomial coefficients and Pascal's identity}
$\binom{n}{j}$ counts the $j$-element subsets of an $n$-element set. The
endpoints satisfy $\binom n0=\binom nn=1$. Pascal's identity is
\[
\binom{k+1}{j}=\binom{k}{j}+\binom{k}{j-1}
\quad(1\le j\le k).
\]
On the Core, these identities are supplied. Your job is to use them with
correct endpoints and indexing.
\end{definitionbox}

\begin{methodbox}{How to manage the endpoints}
For the $(k+1)$-row, separate $j=0$ and $j=k+1$, then apply Pascal's identity
only to $1\le j\le k$:
\[
\sum_{j=0}^{k+1}\binom{k+1}{j}
=1+\sum_{j=1}^{k}\left[\binom{k}{j}+\binom{k}{j-1}\right]+1.
\]
Reindex the second interior sum with $i=j-1$; its range becomes
$0\le i\le k-1$.
\end{methodbox}

\begin{workedbox}{C.1 complete proof: statement and base}
Let
\[
P(n):\quad \sum_{j=0}^{n}\binom nj=2^n,
\qquad n\ge0.
\]
\textbf{Base case.} At $n=0$,
\[
\sum_{j=0}^{0}\binom0j=\binom00=1=2^0,
\]
so $P(0)$ holds.

\textbf{Induction hypothesis.} Fix arbitrary $k\ge0$ and assume
\[
\sum_{j=0}^{k}\binom kj=2^k.
\]
We must prove the corresponding identity for $k+1$.
\end{workedbox}

\columnbreak

\begin{workedbox}{C.1 complete proof: induction step}
\textbf{Step 1: split endpoints and use Pascal.}
\begin{align*}
\sum_{j=0}^{k+1}\binom{k+1}{j}
&=1+\sum_{j=1}^{k}
 \left(\binom kj+\binom{k}{j-1}\right)+1\\
&=1+\sum_{j=1}^{k}\binom kj
 +\sum_{i=0}^{k-1}\binom ki+1\\
&=\sum_{j=0}^{k}\binom kj
 +\sum_{i=0}^{k}\binom ki\\
&=2^k+2^k &&\text{by the IH}\\
&=2^{k+1}.
\end{align*}
The second line uses $i=j-1$ in the second interior sum. The third line
restores the missing endpoint in each copy, so the induction hypothesis can
be used on the fourth line. Thus $P(k+1)$ holds. By mathematical induction,
$\sum_{j=0}^{n}\binom nj=2^n$ for every integer $n\ge0$.
\end{workedbox}

\begin{definitionbox}{Why the reindexing is valid}
When $i=j-1$, the first old index $j=1$ becomes $i=0$, and the last old
index $j=k$ becomes $i=k-1$. The letter itself is a dummy variable; the
limits and summand must change together. Adding the two endpoint $1$s restores
the missing $j=0$ and $j=k$ terms so that two complete $k$-row sums appear.
\end{definitionbox}

\begin{warningbox}{Frequent C errors}
Do not apply Pascal's identity at endpoints where the supplied version does
not apply. Do not change $j$ to $i$ without changing both limits. Do not cite
the IH until the expression is a complete copy of the $k$-case sum. Do not
omit the final induction conclusion.
\end{warningbox}

\begin{checklistbox}{C readiness}
Can I reproduce the proposition, base, IH, endpoint split, reindexing, two
complete sums, explicit IH use, and final conclusion without a missing term?
\end{checklistbox}
\end{multicols}

\clearpage
\section*{More Worked Practice and Whole-Core Readiness}
\begin{multicols}{2}
\begin{workedbox}{Divisibility by induction}
Prove $3\mid(4^n-1)$ for every $n\ge0$.

\textbf{Base.} $4^0-1=0$, divisible by $3$.

\textbf{IH.} Assume $4^k-1=3q$ for some integer $q$. Then
\[
4^{k+1}-1=4(4^k-1)+3=4(3q)+3=3(4q+1).
\]
Because $4q+1$ is an integer, the next expression is divisible by $3$. This
example shows how a definition supplies the usable form of the IH.
\end{workedbox}

\begin{workedbox}{A finite-sum proof}
Prove $1+2+\cdots+n=n(n+1)/2$ for $n\ge1$.

Base: $1=1(2)/2$. Assume
$1+\cdots+k=k(k+1)/2$. Then
\begin{align*}
1+\cdots+k+(k+1)
&=\frac{k(k+1)}2+(k+1)\\
&=\frac{(k+1)(k+2)}2.
\end{align*}
This is the formula with $n=k+1$, so induction completes the proof.
\end{workedbox}

\begin{workedbox}{An inequality with a domain check}
To prove $n!\ge2^{n-1}$ for $n\ge1$, the base is $1=1$. Assuming
$k!\ge2^{k-1}$ for $k\ge1$,
\[
(k+1)!=(k+1)k!\ge2\cdot2^{k-1}=2^k,
\]
because $k+1\ge2$. The domain supplies the comparison needed after the IH.
\end{workedbox}

\columnbreak

\begin{definitionbox}{What the compact Core certifies}
The live Core deliberately has only three problems:
\begin{itemize}
  \item A: a short induction or strong-induction setup;
  \item B: diagnosis and repair of a faulty step;
  \item C: one complete binomial-sum proof.
\end{itemize}
C is essential. Meets requires a fully correct C and one fully correct short
problem. Exceeds requires all three. A retake is a new complete form.
\end{definitionbox}

\begin{methodbox}{Twenty-minute strategy}
\begin{enumerate}
  \item Read all three prompts and write the exact $P(n)$ for C.
  \item Complete A or B cleanly in about 3--4 minutes.
  \item Reserve 12--13 minutes for C.
  \item Use the last minutes to check base values, index limits, where the IH
  was cited, and the final conclusion.
\end{enumerate}
Do not sacrifice the essential C proof to over-polish a short response.
\end{methodbox}

\begin{warningbox}{Proof language that matters}
``Assume true for $n$'' is ambiguous; name the arbitrary index and exact
hypothesis. ``Therefore true'' is incomplete; state that $P(k+1)$ follows and
invoke induction over the stated domain. Algebra alone does not replace the
logical structure.
\end{warningbox}

\begin{checklistbox}{Whole-Core readiness}
I can:
\begin{itemize}
  \item distinguish ordinary from strong induction;
  \item state every component of a valid plan;
  \item find and repair circular use of the target;
  \item calculate the correct next term;
  \item separate endpoints and reindex a binomial sum;
  \item show exactly where the induction hypothesis is used; and
  \item finish with a correctly scoped induction conclusion.
\end{itemize}
\end{checklistbox}
\end{multicols}

\end{document}
